\section{Organization of the Study}
\label{sec:organization_of_the_study}

The study is organized into six chapters.

Chapter~\ref{chap:introduction} introduces the reinforcement-learning exploration problem, motivates the use of intrinsic rewards based on impact and learning progress, states the research objectives, and defines the scope and limitations of the study.

Chapter~\ref{chap:related_literature} reviews related literature on reinforcement learning, PPO, intrinsic motivation, novelty-based exploration, prediction-error curiosity, impact-driven exploration, learning-progress exploration, and noisy dynamics.
It also identifies the research gap addressed by GLPE.

Chapter~\ref{chap:theoretical_framework} presents the theoretical framework.
It formalizes the reinforcement-learning problem, reward augmentation, latent representation learning, forward and inverse dynamics models, feature-space impact, region-local learning progress, the gated exploration criterion, and the assumptions underlying the method.

Chapter~\ref{chap:methodology} describes the methodology.
It details the GLPE method, the PPO training procedure, baseline methods, benchmark environments, evaluation metrics, ablation methodology, project-management considerations, and ethical and reproducibility concerns.

Chapter~\ref{chap:results_and_discussion} presents and discusses the experimental results.
It reports learning curves, step-AUC, final-checkpoint performance, thresholded reliability, ablation results, gating behavior, intrinsic-reward dynamics, and wall-clock efficiency.
The chapter interprets the results in relation to the stated objectives and the design principles of the method.

Chapter~\ref{chap:conclusion_and_recommendations} concludes the study by summarizing the main conclusions, identifying contributions and limitations, and recommending directions for future work.